\begin{beginnerstatement}[In simple terms: Quality and evidence]

Quality assurance comprises the measures used to find errors and check
requirements. A test is a specific check that has been executed. Verification
here means checking in a traceable way whether expected technical behavior is
actually supported by evidence. Evidence is a verifiable result; an artifact
may, for example, be a generated package or a report. A commit is a uniquely
recorded state of the source code, so evidence bound to a commit applies only
to that exact state. \emph{PASS} means passed, \emph{FAIL} means executed and
failed, and \emph{SKIP} means not executed or not applicable. Missing evidence
therefore means \emph{not verified}; it is neither a success nor automatically
an error.

A boundary test deliberately checks values immediately below, at, and above a
boundary, such as 900 and 901 code lines. A fixture is an artificially prepared
test case. If it detects a prohibited import, that specific case is
substantiated; it does not follow that the rule finds every conceivable
architecture violation.

\end{beginnerstatement}
